%% ma: L10-B02-0027
%% lop: 10
%% chuong: 1
%% bai: 2
%% dang: 10.2.2
%% loai: TN4
%% muc: TH
%% dap_an: "C"
%% nguon: "Ảnh giáo viên gửi 10/10/2026 (ThS. Nguyễn Hoàng Việt, Chương 2 Tập hợp), Đề 2 (đề thứ hai, không ghi tiêu đề), Phần 1 Câu 4"
%% kien_thuc: "Bài 2 (tập hợp và các phép toán trên tập hợp)"
%% trang_thai: cho-duyet
%% ly_do: "Dạng toán mới đề xuất 10.2.2: xác định tập hợp bằng cách liệt kê phần tử (giải phương trình, bất phương trình trong N, Z, Q, R)"
\begin{ex}
Hãy liệt kê các phần tử của tập $X=\{x\in\mathbb Q\mid(x^2-x-6)(x^2-5)=0\}$.
\choice{$X=\{-\sqrt5;\sqrt5\}$}{$X=\{-\sqrt5;-2;\sqrt5;3\}$}{\True $X=\{-2;3\}$}{$X=\{\sqrt5;3\}$}
\loigiai{$x^2-x-6=0\Leftrightarrow x\in\{-2;3\}$; $x^2-5=0\Leftrightarrow x=\pm\sqrt5$ là số vô tỉ nên không thuộc $\mathbb Q$. Vậy $X=\{-2;3\}$.}
\end{ex}
